% Insert before \end{document} in supplement_personnel.tex.
% Uses only its existing longtable, booktabs and array packages.
% MR source IDs resolve to the accompanying source_manifest.json and SOURCE_MAP.txt.
\clearpage
\subsection*{C6. Reproducible personnel-profile specifications}
The specifications below define the implementations used for the reported
structural-expression, consensus-response and joint-profile analyses
[MR01--MR03]. For an observation in building
$b$, labels are joined by person and \texttt{heldout\_building}$=b$.
The \texttt{ALL} fit is descriptive only. Class numbers are assigned by the
ordered training-fold class centres, not by matching a test fold to the
\texttt{ALL} solution. Other replacement designs retain the labels specified
for their own comparisons.

Let $S$ denote four raw-response fields:
\begin{sloppypar}
\texttt{turn\_n}, \texttt{short\_turn\_edges},
\texttt{near\_\allowbreak straight\_\allowbreak nodes\_\allowbreak 5deg},
and \texttt{unmatched\_\allowbreak ref\_\allowbreak 9}.
They count floor turns $>5^\circ$, edges shorter than 4\% of the perimeter
with both endpoint turns at least $15^\circ$, turns $\leq5^\circ$, and the
unmatched reference fraction after one-to-one spherical floor-point matching
with a $9^\circ$ acceptance threshold, respectively.
Let $B$ denote \texttt{matched\_\allowbreak top\_\allowbreak bias\_\allowbreak deg},
\texttt{matched\_\allowbreak bottom\_\allowbreak bias\_\allowbreak deg}, \texttt{log\_area\_ratio}, and
\texttt{log\_height}. The signed endpoint deviations use one-to-one periodic
azimuth matching with a $9^\circ$ threshold; the remaining fields are log
submitted/reference floor area and log reconstructed mean height under
camera-height normalization. These are not all absolute-error measures.
\end{sloppypar}

Raw profile features are centred on the median for image, legal reference
and condition. Log time additionally separates origin and time source and
uses positive finite recorded durations only. Each metric is averaged over
valid records within person--target, then equally over targets within image,
images within person--building, and buildings within person. Missing
outcomes are not zero-imputed [MR02].

\begingroup\small
\begin{longtable}{>{\raggedright\arraybackslash}p{.19\linewidth}>{\raggedright\arraybackslash}p{.49\linewidth}>{\raggedright\arraybackslash}p{.23\linewidth}}
\toprule
Profile & Exact feature blocks & Eligibility and class order \\
\midrule\endfirsthead
\toprule Profile & Exact feature blocks & Eligibility and class order \\
\midrule\endhead
Target structural
(\texttt{TARGET\_DSTRUCT})
&
12 features, two domains: overall means of $S$ (4); difficult means and
difficult-minus-easy means of $S$ (8).
One Ward tree has the defined $K=2,3,4$ cuts.
&
Overall $\geq12$ images/3 buildings; easy $\geq5$/2; difficult $\geq4$/2.
Complete required features.
Order by difficult turn count, difficult short-edge count, overall turn count.
\\[3pt]
Consensus response
(\texttt{RFUSE\_GEOM})
&
21 regularized member-effect features: overall, difficult and
difficult-minus-easy coefficients for
\texttt{floor\_turn5\_n},
\texttt{floor\_short\_edge\_03\_n},
\texttt{floor\_redundant5\_n} (structure domain, 9);
\texttt{bottom\_y\_mean\_px}, \texttt{top\_y\_mean\_px},
\texttt{log\_floor\_area}, \texttt{log\_height\_mean}
(boundary/scale domain, 12).
The paper uses the defined $K=2$ cut.
&
Overall $\geq8$ images/3 buildings; easy $\geq5$/2; difficult $\geq3$/2.
Order by difficult fused turn response, then overall bottom-coordinate response.
Native fused-output short edges use the length threshold defined below,
distinct from $S$.
\\[3pt]
Joint profile
(\texttt{OPPJOINT})
&
13 features, four domains: opportunity slopes of $S$ (4), opportunity slopes
of $B$ (4), overall mean/p20 and opportunity slope of peer-centred Q (3),
overall mean and opportunity slope of peer-centred log time (2).
The reported $K=3,4$ are existing cuts.
&
Overall $\geq12$ images/3 buildings; opportunity variance $>10^{-5}$ and
at least two levels; positive-time support $\geq12$/3 with a defined time slope.
Complete required features.
Order by overall Q residual, challenge-turn residual, overall log-time residual.
\\
\bottomrule
\end{longtable}
\endgroup

For valid positive-area fused footprints, the short-edge count is the number
of cyclic edges with Euclidean length strictly below 0.3 camera-height units.
It applies neither a perimeter-normalized threshold nor an endpoint-turn
threshold [MR12].

For the joint profile, the opportunity inputs are
$\log(1+\max(v-4,0))$, $\log(1+H)$ and $\log(1+s)$, where
$v$, $H$ and $s$ are the archived \texttt{ref\_turn5}, \texttt{ref\_H} and
\texttt{ref\_short\_edges} fields.
Training unique image--reference targets receive equal
target-within-image, image-within-building and building weights.
Each axis is centred on its weighted median, divided by its weighted
p85--p15 range (fallback: weighted standard deviation, then 1), clipped to
$[-3,3]$, and the three axes are summed and divided by $\sqrt{3}$.
A person's response slope is weighted covariance of the peer residual and
opportunity divided by weighted opportunity variance. The Q p20 is a weighted
empirical quantile of the same peer residuals, not a raw-Q P10 [MR02].

For each eligible training-fold profile domain, columns are median-centred,
scaled by p85--p15 (fallback: population standard deviation, then 1), and
clipped to $[-3,3]$. Full-SVD PCA retains at most two axes per raw-response
domain and two per consensus-response domain. Each domain is divided by
the training root-mean-squared PC-vector length (floor $10^{-9}$) and by
the square root of the number of domains. Ward linkage uses Euclidean
distances and optimal leaf ordering. The target building is excluded before
history, eligibility, scaling, PCA and clustering. Insufficient histories
remain \texttt{unknown}; they are not assigned to the nearest class.

\paragraph{Regularized consensus-response effects.}
The \texttt{RFUSE\_GEOM} implementation uses ordinary Manual natural
orders, sizes $k\in\{10,12,14,16,20,24\}$, and 24 natural orders per pool.
Repeated deterministic team endpoints are collapsed with their path
occurrence counts retained. A pool--size group requires at least four
distinct eligible teams. The seven geometry outputs and the shared
quality-output eligibility fields must be available; thus Q is not a
GEOM clustering feature, but Q availability affects its training support.
For team $t$ and person $w$, $X_{tw}=1/k_t$ for a member and zero otherwise.
Within pool--size groups, both $X$ and the output matrix $Y$ are centred
using path-occurrence weights. Each building has total weight one, with
equal images within building and equal target--size groups within image;
path occurrences determine weights within each group. The separate
overall/easy/difficult fits are
\[
 \widehat{\mathbf B}
 = (X^\mathsf{T}WX+\lambda I)^{-1}X^\mathsf{T}WY,\qquad
 \lambda=\max\{ \alpha\,\mathrm{tr}(X^\mathsf{T}WX)/P,\;10^{-12}\},
\]
where $P$ is the number of membership columns.
Nested leave-one-building-out selection considers
$\alpha\in\{0.1,1,10\}$ using training-variance-normalized response MSE,
balanced over validation buildings, outputs, strata and domains.
The selection table chooses $\alpha=10$ for the descriptive fit
and all 22 outer folds [MR03]. This is the dimensionless multiplier, not
a fixed numerical penalty of 10. The target-equal
\texttt{TARGET\_RFUSE} and Q-augmented \texttt{RFUSE\_GEOMQ} candidates are
different models and are not substituted here.

\subsection*{C7. Prediction models, preprocessing and test support}
The following rows identify the exact implementations behind the reported
comparisons. A shared label such as \emph{history} does not imply identical
aggregation across these analyses [MR04--MR07].

\begingroup\small
\begin{longtable}{>{\raggedright\arraybackslash}p{.19\linewidth}>{\raggedright\arraybackslash}p{.49\linewidth}>{\raggedright\arraybackslash}p{.23\linewidth}}
\toprule
Comparison & Features, estimator and preprocessing & Split, support and estimand \\
\midrule\endfirsthead
\toprule Comparison & Features, estimator and preprocessing & Split, support and estimand \\
\midrule\endhead
Difficulty versus computable geometry
&
Weighted linear ridge, penalty 1 on every coefficient except the intercept.
Controls: training-level person and record-source indicators; medium/hard
difficulty indicators. Added numeric features:
\texttt{ref\_turn5}, \texttt{ref\_H}, \texttt{ref\_short\_edges},
\texttt{ref\_min\_horizontal\_deg}, short-edge presence, and minimum
short-edge span conditional on presence (zero for structural absence).
Numeric weighted mean/SD scaling is training-only; no tuning.
The binary simplification prediction is clipped linear regression,
not logistic regression.
&
22 building-held-out folds; 6,512 records, 206 images, 60 people.
Training images have equal total weight; test errors are averaged within
image then building, with buildings equal.
MSE for turn deficit and simplification; MAE for Q.
Missing outcomes are excluded per outcome; nonfinite geometry is excluded.
Unseen people have zero person indicators.
\\[3pt]
Same-room image-mean Q:
9.496 to 6.464
&
Direct source-image mean Q substitutes for the target-image mean.
The baseline is the mean of other-building image means.
No fitted regression coefficient, scaling, imputation or hyperparameter.
&
128 directly adopted, non-pending pairs, both directions:
256 predictions, 65 images, 13 buildings.
Average absolute errors over source directions per target,
then images within building and buildings equally.
\\[3pt]
Same-room individual Q:
15.645 to 13.385
&
Baseline: the person's record-mean Q outside the target building.
Let $\theta_w$ be the record-mean outside-building residual from
leave-one-person-out task-peer Q.
Prediction: source-image mean Q of other people plus $\theta_w$,
clipped to $[0,100]$.
This is \texttt{public\_external} with a single direct source,
not the half-own-residual variant.
&
1,793 target answers, 65 images, 56 people, 13 buildings;
5,198 source-specific prediction rows.
Shared valid support across the saved predictors; missingness is not zero.
Source-specific errors are averaged per target answer before
image/building aggregation.
\\[3pt]
Same-room structure:
short edges 0.526 to 0.472
&
For each target answer, average the person's feature over directly
neighbouring source views on which that person answered.
The comparator averages other people's features on the same source views.
No fitted parameters. This averages source features before evaluating
the target error, unlike the single-direct Q comparison.
&
Same 1,793/65/56/13 support, requiring outside-building personal history.
The signed floor-offset comparison is in ERP pixels and uses spherical
bottom-point matching; it is not the azimuth-paired degree feature in $B$.
\\[3pt]
Recorded purpose added to geometry:
5.158 versus 5.379
&
Five numeric features: reference turn count, log(1+\texttt{ref\_H}),
short-edge/turn fraction, log minimum horizontal span and log
90th-percentile floor radius (last two floored at $10^{-5}$ before log).
Difficulty is categorical; the augmented model adds purpose.
Training-median imputation with missing indicators, standard scaling,
one-hot categories with unseen levels ignored.
Fixed candidate portfolio detailed below.
&
Outer LOBO on 206 ordinary images/22 buildings; reported repeated-purpose
subset: 190 images/21 buildings.
Inner building-grouped five-fold model selection.
Q predictions clipped to $[0,100]$; reported MAE is image-equal.
Geometry is reference-derived, not a no-reference RGB predictor.
\\
\bottomrule
\end{longtable}
\endgroup

The recorded-purpose portfolio contains linear ridge
$\alpha\in\{0.1,1,10,100\}$, quadratic-numeric ridge
$\alpha\in\{10,100\}$, and histogram gradient boosting
(120 iterations, 7 leaf nodes, minimum 15 samples per leaf,
learning rate 0.05, L2 penalty 10, seed 2026102703).
Quadratic expansion follows imputation/scaling and has degree two with no
bias column. The inner criterion is the mean of held-out-fold MAEs.
Every preprocessing step is re-fitted inside each inner training fold.
For Q, the geometry-only model selected quadratic ridge $\alpha=10$ in
16 outer folds and boosting in 6; geometry plus purpose selected quadratic
ridge $\alpha=10$ in 18, boosting in 3, and linear ridge $\alpha=1$ in 1.
The archived 44-row model table identifies every fold [MR07].

\subsection*{C8. Continuous Q and Q/T selection}
The fixed candidate roster is shared by the policies and does not expand
when the historical selection criteria change.
All target-building records are excluded from historical estimates.
For each person, $q_w$ is the record-mean Q residual relative to other people
on the same image, condition, legal target and reference version.
The cost $c_w$ is the exponential of the record-mean log-time residual,
additionally separating origin and time source. Q and time histories
each require at least eight images and three buildings.
The main panel uses the compatible complete-time history subset, without
mixing incompatible time sources [MR08].

Q-only selects the $n$ largest $q_w$, breaking ties by person identifier.
Q/T starts from that team and repeatedly chooses the best strictly
cost-reducing one-person exchange subject to a fixed historical mean-Q
floor of the initial team's mean minus 2 points. Feasible exchanges are
ordered by total $c_w$, then descending mean $q_w$, then sorted person IDs;
cost must decrease by more than $10^{-9}$. It stops when no such exchange
exists. This is a deterministic local-search rule, not a globally optimal
solution, and target Q does not select the team. Each
$n\in\{3,4,5,6,7,9,13\}$ is optimized separately, with more than $n$
eligible candidates; teams across sizes need not be nested.

The primary common-support identity is
(source--target pair, source pool, target pool, domain).
Both policies must have been executed at every listed size, with exactly
the same candidate roster throughout. Only ordinary Manual, comparable,
directly adopted relations without pending or substantive scope conflicts
are included. This yields 126 contexts, 46 target images and 9 buildings.
Let $\mathcal A$ average contexts within target image, images within
building, and buildings equally. The two primary estimands are
\[
 \mathcal A(Q_{\rm QT}-Q_{\rm Q}),\qquad
 \exp\{\mathcal A[\log(T_{\rm QT}/T_{\rm Q})]\}.
\]
Thus the labor ratio is not a ratio of arithmetic mean labor.
At $n=5$, these are $0.21836$ and $0.81343$, with empirical
building-cluster percentile intervals $[-0.53836,0.86896]$ and
$[0.76045,0.88192]$. The archive uses 20,000 building-bootstrap
resamples, seed 20261030, synchronized across sizes [MR09].
These quantify empirical cluster sensitivity, not a guarantee of
population-sampling coverage.

The $n=5$ versus $n=13$ comparison additionally requires both policies to
be scoreable at every size: 124 contexts, the same 46 images and 9 buildings.
Its labor-growth estimand is
$\exp\{\mathcal A[\log(T_{13}/T_5)]\}$, separately for each policy
(2.33142 and 2.17016). It compares independently selected teams,
not the causal effect of adding eight people to a fixed team.
Unscoreable Q is never replaced by zero.

\subsection*{C9. Source-history stopping and acquisition cost}
The reported 86.054 versus 85.189 Q comparison is the archived
\texttt{maximum\_history}, \texttt{whole\_group},
\texttt{all\_attempts}, non-pending-direct-relations panel.
For a room group containing $v$ views, it uses the maximum available
source-view history $m=v-1$, not one source view [MR10].
For each shared member-order replay $r$, let
\[
 s_r(k)=\frac{1}{m}\sum_{h\in H} Q_{hr}(k),\qquad
 M_r=\max_{j\geq3:\,s_r(j)\ {\rm finite}}s_r(j).
\]
\[
 k_{\rm req,r}=
 \min\{k\geq3:s_r(k)\ {\rm finite},\ s_r(k)\geq M_r-2\}.
\]
All source values at a prefix must be available; the implementation uses
ordinary means, no interpolation, and a numerical comparison tolerance
of $10^{-10}$. If no valid prefix exists, no choice is produced.
The target then advances from $k_{\rm req,r}$ to its first complete
geometric output: status \texttt{ok}, at least four nodes, and an available
connection. A target with no subsequent complete output remains unavailable.
Target reference Q is not used by this decision. The fixed-seven comparator
requests $\min(7,K)$ and applies the same geometric top-up, where $K$ is
the available common-worker pool size.
The separate online rule triggers at $k\geq5$ when three consecutive
complete outputs have equal node counts; otherwise it uses the full pool.
Equal counts do not establish equal geometry.

The quality panel contains five same-room groups, 21 images and five
buildings (5,376 replay attempts). Outcomes are first summarized
within each target view over its replay/history choices, then equally
over target views within group, groups within building, and buildings.
The compatible complete-time panel contains only three groups,
11 images and three buildings. Let $L_{vr}(k)$ be the recorded member
labor through prefix $k$, and $k_{7,vr}$ the fixed-seven-plus-geometry choice.
For a source/target split $(H,T)$, the project acquisition ratio is
\[
 \rho_r=
 \frac{\sum_{h\in H}L_{hr}(K)+
       \sum_{t\in T}L_{tr}(k_{{\rm stop},tr})}
      {\sum_{h\in H}L_{hr}(k_{7,hr})+
       \sum_{t\in T}L_{tr}(k_{7,tr})}.
\]
The archive averages these ratios over replays, splits, groups and buildings,
yielding 2.11526. This is an arithmetic mean of project ratios,
not a ratio of the target-only mean labor values
(1,660.3 versus 1,307.9 person-seconds), and not the geometric ratio in C8.
Replays reuse collected annotations and are not charged as repeated
real acquisitions. Unmeasured reference-production and assisted-connection
costs are not imputed as zero; recorded person-time is not wall-clock delay.

\paragraph{Source and implementation record.}
The source author is Sparkling-Flames (2026), \emph{3D\_Manhattan\_label},
a research software and output repository. All MR source IDs above resolve
to the full file titles, paths, Git blob identifiers and fixed-commit URLs
in the accompanying \texttt{source\_manifest.json} and
\texttt{SOURCE\_MAP.txt}, at commit
\texttt{a116cebbd03321f55c757130a4dd88b299539155}.
MR01--MR03 identify the profile definitions, result bindings and
underlying profile code; MR04--MR07 identify prediction scripts and matched
output tables; MR08--MR10 identify selection/stopping implementations,
common-support tables and cost summaries.
The source map links each method specification to its implementation and
reported output tables.

